\section{Cinemática Diferencial e Jacobiano}
\label{sec_cinematica_diferencial_jacobiano}

% A cinemática diferencial relaciona as velocidades articulares
% $\dot{\vec q} = [\dot\theta_1,\dot\theta_2,\dot\theta_3]^T$
% com a velocidade linear da ferramenta no referencial da espaçonave $\{B\}\equiv\{0\}$, coerente com a cinemática direta implementada.

% \subsection{Jacobiana linear e angular}
O mapeamento das velocidades articulares é dado pela Jacobiana linear $\mathbf{J}_{v}(\vec q)$:

\begin{equation}
\label{eq_jacob_linear1}
{}^{0}\!\dot{\vec{x}}_{E} = \mathbf{J}_{v}(\vec q)\,\dot{\vec q}.
\end{equation}

Na formulação geral, a Jacobiana geométrica reúne ambas as contribuições (linear e angular) na mesma equação:

\begin{equation}
\begin{bmatrix}
{}^{0}\!\dot{\vec{x}}_{E} \\
{}^{0}\!\vec{\omega}_{E}
\end{bmatrix}
=
\begin{bmatrix}
\mathbf{J}_{v}(\vec q) \\
\mathbf{J}_{\omega}(\vec q)
\end{bmatrix}
\dot{\vec q}.
\end{equation}

Em que $\mathbf{J}_{\omega}(\vec q)\in\mathbb{R}^{3\times 3}$ é o Jacobiano angular.
Para juntas revolutas, cada coluna de $\mathbf{J}_{\omega}$ coincide com o eixo unitário de rotação da junta correspondente, expresso no referencial da espaçonave $\{B\}\equiv\{0\}$.

\subsection{Construção do Jacobiano}

As colunas de $\mathbf{J}_{v}$ são obtidas a partir dos eixos unitários de cada junta, expressos no referencial da espaçonave $\{B\}\equiv\{0\}$, e das posições dos respectivos pontos de rotação.
Denotado por ${}^{0}\!\hat{u}_{i}$, o eixo da junta $i$ expresso em $\{0\}$ e por ${}^{0}\!\vec p_{i}$ a posição do ponto sobre esse eixo, extraídos da cinemática direta, a contribuição linear de uma junta revoluta $i$ é dada por:

\begin{equation}
J_{v,i}
=
{}^{0}\!\hat{u}_{i}
\times
\bigl(
{}^{0}\!\vec p_{E}
-
{}^{0}\!\vec p_{i}
\bigr).
\end{equation}

O termo ${}^{0}\!\vec p_{E}$ é a posição do efetuador no mesmo referencial.
No manipulador 3R deste trabalho, a junta 1 rotaciona em torno de um eixo vertical alinhado com $\hat{z}_{B}$, enquanto as juntas~2 e~3 rotacionam em torno de eixos paralelos a $\hat{x}_{B}$. 
Portanto, cada coluna de $\mathbf{J}_{v}$ representa a velocidade linear instantânea do efetuador decorrente de uma velocidade unitária aplicada em torno desses eixos físicos de rotação.

\section{Interpretação geométrica da Jacobiana}
\label{sec_interpretacao_geometrica_jacobiana}

A Jacobiana descreve como as velocidades articulares geram velocidade linear e angular no efetuador, embora, neste trabalho, a interpretação geométrica seja enfatizada principalmente na componente linear.
Cada coluna possui interpretação geométrica direta, conforme \cite{craig2005}.

\subsection{Contribuição da junta revoluta}

Para uma junta revoluta, a rotação ocorre em torno do eixo ${}^{0}\!\hat{z}_{\,i-1}$ do referencial anterior, que na convenção de Denavit--Hartenberg coincide com o eixo da junta.
Essa rotação produz uma velocidade linear obtida pelo produto vetorial entre esse eixo e o vetor que conecta o eixo ao efetuador:

\begin{equation}
J_{v,i}
=
{}^{0}\!\hat{z}_{\,i-1}
\times
\left(
{}^{0}\!\vec{p}_{E}
-
{}^{0}\!\vec{p}_{\,i-1}
\right).
\end{equation}

O termo $J_{v,i}$ representa o movimento linear imposto ao efetuador, cujo efeito cresce com a distância perpendicular ao eixo ${}^{0}\!\hat{z}_{\,i-1}$.

\subsection{Síntese geométrica}

Cada coluna da Jacobiana linear representa a velocidade instantânea do efetuador causada por uma velocidade unitária na junta $i$:

\begin{equation}
\dot{q}_{i} = 1,
\qquad
\dot{q}_{j} = 0,\; j \neq i.
\end{equation}

\subsection{Jacobiano angular}

De forma análoga, a componente angular da Jacobiana, $\mathbf{J}_{\omega}(\vec q)$, relaciona as velocidades articulares à velocidade angular do efetuador ${}^{0}\!\vec{\omega}_{E}$. Para uma junta revoluta $i$, a contribuição angular é dada diretamente pelo eixo de rotação expresso no referencial da espaçonave:
\begin{equation}
J_{\omega,i} = {}^{0}\!\hat{u}_{i}.
\end{equation}

O termo ${}^{0}\!\hat{z}_{i}$ é o eixo unitário da junta $i$. No manipulador 3R, obtém-se:

\begin{equation}
\mathbf{J}_{\omega}(\vec q)=
\bigl[\,{}^{0}\!\hat{z}_{1}\;\;{}^{0}\!\hat{z}_{2}\;\;{}^{0}\!\hat{z}_{3}\,\bigr].
\end{equation}

Sendo ${}^{0}\!\hat{z}_{1}$ alinhado a $\hat{z}_{B}$ e ${}^{0}\!\hat{z}_{2}$, ${}^{0}\!\hat{z}_{3}$, normal a convenção \gls{dh}, a $\hat{x}_{B}$.
Essa estrutura reflete diretamente a distribuição física dos eixos das juntas do \gls{mr} acoplado à espaçonave.

\section{Singularidades e interpretação geométrica}
\label{sec_singularidades_jacobiana}

Singularidades ocorrem quando a Jacobiana linear perde posto, tornando impossível gerar determinadas velocidades cartesianas.
Ocorrem quando eixos de juntas tornam-se alinhados ou quando o efetuador atinge uma configuração de liberdade reduzida.

Quando colunas de $\mathbf{J}_{v}$ tornam-se linearmente dependentes, o espaço de velocidades possíveis se reduz. Pequenas movimentações desejadas podem exigir velocidades articulares elevadas. Conforme \cite{craig2005}, a independência direcional dos eixos das juntas é essencial para manter posto pleno.

A partir dessa relação direta, surge o problema inverso:
determinar $\dot{\vec{s}}$ que produza uma velocidade cartesiana desejada ${}^{0}\!\dot{\vec{x}}_{E,des}$.
Idealmente, teríamos
$\dot{\vec{s}} = \mathbf{J}_{v}^{-1}(\vec{s})\,{}^{0}\!\dot{\vec{x}}_{E,des}$,
mas, em geral, $\mathbf{J}_{v}$ não é quadrada ou pode perder posto.
Adota-se então a pseudo-inversa amortecida \cite{sciavicco_siciliano2010}:

\begin{equation}
\label{eq_pseudoinversa_preIK}
\dot{\vec{s}}
=
\mathbf{J}_{v}^{T}
\bigl(
\mathbf{J}_{v}\mathbf{J}_{v}^{T}
+
\lambda^{2}\mathbf{J}_{3}
\bigr)^{-1}
\,{}^{0}\!\dot{\vec{x}}_{E,des}.
\end{equation}

Sendo $\lambda > 0$ atuando como termo de regularização perto de singularidades e $\mathbf{J}_{3}$ denotando a matriz identidade de ordem três.


Essa formulação permite projetar a velocidade desejada no espaço articular de forma estável, estabelecendo a ponte entre cinemática diferencial e cinemática inversa diferencial.
